\documentclass[11pt,a4paper]{article}
\usepackage[margin=2.5cm]{geometry}
\usepackage{amsmath,amssymb,bm}
\usepackage{booktabs}
\usepackage{xcolor}
\usepackage[colorlinks=true,linkcolor=blue!50!black,urlcolor=blue!60!black,citecolor=blue!50!black]{hyperref}
\usepackage{enumitem}
\setlist{nosep,leftmargin=*}

\newcommand{\QP}{Q_{\mathrm{P}}}
\newcommand{\aeff}{a}
\newcommand{\code}[1]{\texttt{#1}}

\title{Monte-Carlo statistical uncertainty in the SolarOscFitter likelihood\\[2pt]
\large The Dembinski--Abdelmotteleb and Conway (Barlow--Beeston) treatments, and how they are implemented}
\author{Notes for the SolarOscFitter analysis}
\date{October 2026}

\begin{document}
\sloppy
\maketitle

\begin{abstract}
The predicted event counts in the fit are built from finite Monte-Carlo (MC) samples: signal
response matrices and background templates. Their statistical fluctuations are not in a plain
Poisson likelihood, which treats the prediction as exact. This note explains the method of
Dembinski and Abdelmotteleb~\cite{DA}, which derives a fast approximate likelihood from the
exact one of Barlow and Beeston~\cite{BB}, together with Conway's earlier
approximation~\cite{Conway}. It then describes how the fitter implements both
(\code{mc\_stat\_likelihood}), including the adaptation needed because the day bin and all night
bins of one reconstructed-energy bin are built from the same MC events.
\textbf{The paper:} H.~Dembinski and A.~Abdelmotteleb, \emph{A fast and stable approximate
maximum-likelihood method for template fits},
\href{https://arxiv.org/abs/2206.12346}{arXiv:2206.12346}.
\end{abstract}

\section{The problem}

In each analysis bin the fit compares an observed count $n$ with a prediction $\mu(\bm p)$ that
depends on the parameters $\bm p$ (oscillation parameters, fluxes, nuisances). The plain Poisson
likelihood assumes $\mu$ is known exactly. But $\mu$ is computed from simulation: a response
matrix filled with a finite number of MC events, background templates filled with weighted MC
events. If the MC sample were generated again, $\mu$ would change. When that MC fluctuation is
not small compared with the data fluctuation $\sqrt{n}$, ignoring it makes the fit too confident:
contours shrink and pulls are too wide. With Asimov data the best fit does not move, but the
sensitivity is overstated.

The standard remedy treats the true template content in each bin as unknown, constrained by the
MC that was actually generated. The question is how to do that without introducing a free
parameter for every template in every bin.

\section{Baker--Cousins form}

All likelihoods below are written, following Baker and Cousins~\cite{BakerCousins}, as
\begin{equation}
  Q(\bm p) = -2\ln\frac{L(n;\mu(\bm p))}{L(n;n)},
\end{equation}
whose minimum is asymptotically $\chi^2$-distributed. For one Poisson bin,
\begin{equation}
  \QP(n;\mu) = 2\left[\mu - n + n\ln\frac{n}{\mu}\right] ,
\end{equation}
with the convention $n\ln(n/\mu)=0$ for $n=0$. This is exactly what the fitter's
\code{poissonLogLikelihood} computes (as $-\tfrac12\QP$).

\section{The exact Barlow--Beeston likelihood}

Let the prediction in a bin be a sum of components (templates) $k$ with yields $y_k$. Template $k$
has an observed MC count $a_k$ in the bin, out of $M_k$ in total, and an unknown true expected MC
count $\xi_k$. Then $\mu = \sum_k y_k\,\xi_k/M_k$, and the MC counts are themselves Poisson
around $\xi_k$. The exact likelihood of data and simulation together is~\cite{BB,DA}
\begin{equation}
  Q = \QP(n;\mu) + \sum_k \QP(a_k;\xi_k).
  \label{eq:bb}
\end{equation}
Every $\xi_k$ in every bin is a nuisance parameter. Barlow and Beeston profile them by solving a
non-linear equation per bin. This is exact, but slow when there are many bins, and the finite
accuracy of the solver can make the likelihood slightly discontinuous, which upsets gradient-based
minimisers and samplers.

\section{Conway's approximation (``Barlow--Beeston lite'')}

Write $\xi_k = \beta_k a_k$. Conway's central approximation~\cite{Conway} is to use one common
factor $\beta_k\approx\beta$ per bin:
\begin{equation}
  \mu = \sum_k \frac{y_k\beta_k a_k}{M_k} \approx \beta \underbrace{\sum_k \frac{y_k a_k}{M_k}}_{\mu_0}
  = \beta\mu_0 .
\end{equation}
This holds when one component dominates the bin. Inserting it in Eq.~\eqref{eq:bb} gives
(\cite{DA}, Eq.~4)
\begin{equation}
  Q \approx \QP(n;\beta\mu_0) + 2a\,\bigl[(\beta-1) - \ln\beta\bigr], \qquad a=\sum_k a_k .
  \label{eq:branch}
\end{equation}
Conway then expands $\ln\beta$ to second order around $\beta=1$, so the constraint becomes Gaussian:
\begin{equation}
  Q = \QP(n;\beta\mu_0) + \frac{(\beta-1)^2}{V(\beta)}, \qquad V(\beta)\approx\frac1a
  = \frac{\sum w^2}{\mu_0^2}.
  \label{eq:conway}
\end{equation}
$V(\beta)$ is the squared relative MC uncertainty of the bin; with weighted MC it is
$\sum_i w_i^2/(\sum_i w_i)^2$. Setting $\partial Q/\partial\beta=0$ gives a quadratic with one
valid root, so $\beta$ is profiled in closed form. This is HistFactory's \code{staterror}, and
MaCh3's \code{kBarlowBeeston}.

The Gaussian expansion is good when $a$ is large, but frequently the MC sample is \emph{not} large
in every bin: then the constraint is too weak in the direction of small $\beta$ and too strong
for large $\beta$.

\section{The Dembinski--Abdelmotteleb approximation}

Dembinski and Abdelmotteleb~\cite{DA} start from the same Eq.~\eqref{eq:branch}, but do
\emph{not} expand the logarithm. Minimising over $\beta$ directly,
\begin{equation}
  \frac{\partial Q}{\partial\beta}=0 \;\Rightarrow\;
  \mu_0 - \frac{n}{\beta} + a\Bigl(1-\frac1\beta\Bigr) = 0
  \;\Rightarrow\;
  \boxed{\beta = \frac{n+a}{\mu_0+a}} ,
  \label{eq:da-beta}
\end{equation}
which is even simpler than Conway's root (\cite{DA}, Eq.~11). The interpretation is direct:
\begin{itemize}
  \item if $a\ll n$ (little MC), $\beta$ simply moves $\mu_0$ to $n$, and the bin says nothing about
        the parameters;
  \item if $a\to\infty$ (perfect MC), $\beta\to1$ and $Q\to\QP(n;\mu_0)$, the plain Poisson likelihood.
\end{itemize}
With a few lines of algebra, Eq.~\eqref{eq:branch} can be rewritten as (\cite{DA}, Eq.~12)
\begin{equation}
  \boxed{Q = \QP(n;\beta\mu_0) + \QP(a;\beta a)} .
  \label{eq:da}
\end{equation}
This form is the paper's key point: \textbf{data and simulation are treated symmetrically, both as
Poisson counts}. The MC enters as if the template were an independent measurement of $a$ events
whose expectation is $\beta a$. Poisson-based likelihoods are known to behave better than Gaussian
ones for small counts, so this likelihood is expected to outperform Conway's when the MC is small,
at the same computational cost.

\subsection*{Weighted MC and the effective count}

MC events usually carry weights. Following Bohm and Zech~\cite{BohmZech}, a sum of weights
$\sum_i w_i$ is treated as a scaled Poisson variable with \emph{effective count}
\begin{equation}
  n_{\mathrm{eff}} = \frac{\bigl(\sum_i w_i\bigr)^2}{\sum_i w_i^2}.
\end{equation}
The paper's improved version (\cite{DA}, Sec.~5) treats the whole prediction $\mu_0$ as one
template whose MC variance is $V_\mu$, with effective count
\begin{equation}
  a = \mu_{\mathrm{eff}} = \frac{\mu_0^2}{V_\mu}, \qquad
  Q = \QP(n;\beta\mu_0) + \QP(\mu_{\mathrm{eff}};\beta\mu_{\mathrm{eff}}), \qquad
  \beta = \frac{n+\mu_{\mathrm{eff}}}{\mu_0+\mu_{\mathrm{eff}}}
  \label{eq:da-improved}
\end{equation}
(\cite{DA}, Eqs.~17--18, with unweighted data). This removes the ``one dominant component''
restriction, and it is the form implemented here and in MaCh3.

\subsection*{What the paper finds}

The authors fit two toy problems with exact Barlow--Beeston (ROOT's \code{TFractionFitter}),
Conway, and their new likelihood, for MC samples of 100, 1000 and 10\,000 events per template,
and measure the bias and the coverage (pull width) of the fitted yields. Their conclusions:
\begin{itemize}
  \item with large MC samples all three agree, because the MC uncertainty becomes negligible;
  \item with small, Poisson-distributed templates the approximate likelihoods perform comparably to
        the exact one; Conway's produces the most outliers;
  \item with \emph{weighted} templates the exact Barlow--Beeston likelihood underestimates the
        variance of the fitted yields, while both generalised approximate likelihoods match it;
  \item the approximations are much faster than exact Barlow--Beeston once there are more than
        about 200 bins, and the new one is about 25\,\% faster than Conway's for few bins.
\end{itemize}

\section{How MaCh3 does it}

MaCh3 (\href{https://github.com/mach3-software/MaCh3}{github.com/mach3-software/MaCh3}) offers
the test statistic as a configuration choice (\code{LikelihoodOptions: TestStatistic}):
\code{Poisson}, \code{BarlowBeeston} (Conway), \code{DembinskiAbdelmotteleb}, \code{IceCube}
(the marginal likelihood of Arg\"uelles et al.~\cite{IceCube}) and \code{Pearson}. All are
evaluated per analysis bin in \code{SampleHandlerInterface::}\allowbreak\code{GetTestStatLLH(data, mc, w2)}, where
\code{w2} is the sum of the squared weights of all MC events (signal and background) filled into
the bin. By default \code{w2} is computed once and then frozen (\code{UpdateW2 = false}). The
fitter's per-bin (``independent'') option reproduces MaCh3's formulas exactly; this is checked
in \code{tests/test\_mcstat.jl}.

\section{Implementation in SolarOscFitter}

\subsection{Effective MC counts}

For each channel and reconstructed-energy bin $j$, the prediction at the Asimov truth, summed over
day and all night bins, $M_j$, is rebuilt from its components, and its MC variance $V_j$ is
computed (\code{mc\_effective\_counts} in \code{src/likelihoods/mcstat.jl}). The effective count is
$a_j = M_j^2/V_j$, computed once at the truth and then held fixed, as in MaCh3.
\begin{description}
  \item[Backgrounds.] Each template's bin content is a sum of selected event weights, so its
    relative variance is $\sum w^2/(\sum w)^2$; it is multiplied by the square of the template's
    contribution at its true normalisation.
  \item[Signal.] The response $R_{ij} = n_{ij}/N_i$ is filled with $n_{ij}$ MC events and normalised
    per true-energy row ($N_i=\sum_j n_{ij}$). Given $N_i$ a row is multinomial, so
    $\mathrm{Var}(n_{ij}) = n_{ij}(1-n_{ij}/N_i)$. An MC event of row $i$ stands for $T_i/N_i$
    predicted events, where $T_i$ is the oscillated true-energy rate (day + night), so
    \begin{equation}
      V_j^{\mathrm{resp}} = \sum_i n_{ij}\Bigl(1-\frac{n_{ij}}{N_i}\Bigr)\Bigl(\frac{T_i}{N_i}\Bigr)^2
      (\varepsilon_j g_j)^2 ,
    \end{equation}
    with $\varepsilon_j$ the selection efficiency and $g_j$ a fixed per-bin factor (forward
    fraction with the scattering-angle cut, semi-inclusive split). The efficiency, a ratio of MC
    counts $\varepsilon_j = s_j/t_j$, adds the binomial term $c_j^2(1-\varepsilon_j)/(\varepsilon_j t_j)$.
  \item[Modes.] Plain, angular, inclusive and semi-inclusive channels are covered; the
    semi-inclusive mis-identification term enters $M_j$ without MC variance.
\end{description}
The components are checked to reproduce the Asimov data bin by bin, in every mode
(\code{tests/derived\_mode.jl}), and the signal variance is checked against brute-force resampling
of the MC (\code{tests/test\_mcstat.jl}).

\subsection{One factor per energy bin for day and night}

The night prediction in each zenith bin is the same set of MC events as the day prediction,
weighted by exposure. Their MC fluctuations are therefore not independent: if the MC happened to
put too many events into energy bin $j$, it did so for the day bin and for every night bin. Giving
each $(\cos z, E)$ bin its own factor, as a literal per-bin treatment does, counts some 50 copies
of one fluctuation as independent, and so understates the effect. The default
(\code{mc\_stat\_day\_night: shared}) therefore uses one $\beta_j$ for the group $G_j$ of all day
and night bins of energy bin $j$:
\begin{equation}
  Q_j = \sum_{b\in G_j}\QP(n_b;\beta_j\mu_b) + \QP(a_j;\beta_j a_j)
  \quad\text{(Dembinski--Abdelmotteleb)}.
\end{equation}
Because $\partial\QP(n;\beta\mu)/\partial\beta = 2(\mu - n/\beta)$ is linear in $1/\beta$, the
minimisation stays in closed form:
\begin{equation}
  \sum_b\Bigl(\mu_b - \frac{n_b}{\beta}\Bigr) + a\Bigl(1-\frac1\beta\Bigr) = 0
  \;\Rightarrow\;
  \beta_j = \frac{N_j + a_j}{M_j + a_j}, \qquad N_j=\sum_{b\in G_j} n_b,\; M_j = \sum_{b\in G_j}\mu_b .
\end{equation}
For Conway's form, $Q_j = \sum_b\QP(n_b;\beta\mu_b) + a_j(\beta-1)^2$, the same step gives the
quadratic $a_j\beta^2 + (M_j-a_j)\beta - N_j = 0$. With $\beta = 1+x$ both are evaluated in a
form that does not lose precision when $a_j$ is large:
\begin{align}
  \text{DA:}&\quad x = \frac{N_j - M_j}{M_j + a_j}, &\tfrac12\QP(a;\beta a) &= a_j\,\bigl(x - \ln(1+x)\bigr),\\
  \text{Conway:}&\quad x = \frac{-2(M_j-N_j)}{(a_j+M_j) + \sqrt{(a_j-M_j)^2 + 4a_jN_j}}, &
  \tfrac12 a_j(\beta-1)^2 &= \tfrac12 a_j x^2 .
\end{align}
No iterative solver is needed. The likelihood remains smooth and is differentiated exactly by
ForwardDiff and by Enzyme.

\subsection{Properties checked by the tests}

\begin{itemize}
  \item the closed-form $\beta$ is the profile minimum (against brute-force minimisation over
        $\beta$), for single bins and for shared groups;
  \item the per-bin form equals MaCh3's \code{GetTestStatLLH} for both methods;
  \item $a\to\infty$ recovers the plain Poisson likelihood; with Asimov data $\beta=1$ and the
        likelihood is unchanged at the truth;
  \item profiling can only raise the likelihood relative to Poisson, and less so for larger $a$;
  \item the per-bin diagnostic pieces add up to the channel likelihood;
  \item the full likelihood with each option is pinned by fixture regressions, and Enzyme agrees
        with ForwardDiff.
\end{itemize}

\subsection{Usage}

\begin{center}\small
\begin{tabular}{lll}
\toprule
key & values & meaning \\
\midrule
\code{mc\_stat\_likelihood} & \code{Poisson} (default) & no MC uncertainty \\
 & \code{DembinskiAbdelmotteleb} & Eq.~\eqref{eq:da-improved} (recommended) \\
 & \code{BarlowBeeston} & Conway, Eq.~\eqref{eq:conway} \\
\code{mc\_stat\_day\_night} & \code{shared} (default) & one $\beta$ per energy bin, day + night \\
 & \code{independent} & one $\beta$ per bin (MaCh3's per-bin treatment) \\
\bottomrule
\end{tabular}
\end{center}
Both keys can be set at the top level of the configuration (default for every detector) or in a
detector block. The setup log reports the median and maximum relative MC uncertainty
$1/\sqrt{a_j}$ in the fit window, and the per-bin breakdown CSV lists $a_j$.

\subsection{Limitations}

\begin{itemize}
  \item The effective counts are frozen at the truth. Away from the truth, the MC variance of a
        reweighted prediction changes; MaCh3 makes the same default choice.
  \item The signal efficiency's MC uncertainty is treated as independent of the response's; the
        semi-inclusive mis-identification response and the scattering-angle response carry no MC
        variance.
  \item Pre-binned input files (histograms instead of events) have no meaningful
        $\sum w^2$: each row looks like a single heavy event, which makes their MC uncertainty
        appear huge. The inputs must be event-level.
\end{itemize}

\begin{thebibliography}{9}
\bibitem{DA} H.~Dembinski and A.~Abdelmotteleb, \emph{A fast and stable approximate
  maximum-likelihood method for template fits},
  \href{https://arxiv.org/abs/2206.12346}{arXiv:2206.12346} [stat.ME] (2022).
\bibitem{BB} R.~Barlow and C.~Beeston, \emph{Fitting using finite Monte Carlo samples},
  Comput.\ Phys.\ Commun.\ \textbf{77} (1993) 219.
\bibitem{Conway} J.~S.~Conway, \emph{Incorporating nuisance parameters in likelihoods for
  multisource spectra}, PHYSTAT 2011, \href{https://arxiv.org/abs/1103.0354}{arXiv:1103.0354}.
\bibitem{BakerCousins} S.~Baker and R.~D.~Cousins, \emph{Clarification of the use of chi-square and
  likelihood functions in fits to histograms}, Nucl.\ Instrum.\ Meth.\ \textbf{221} (1984) 437.
\bibitem{BohmZech} G.~Bohm and G.~Zech, \emph{Statistics of weighted Poisson events and its
  applications}, Nucl.\ Instrum.\ Meth.\ A \textbf{748} (2014) 1.
\bibitem{IceCube} C.~A.~Arg\"uelles, A.~Schneider and T.~Yuan, \emph{A binned likelihood for
  stochastic models}, JHEP \textbf{06} (2019) 030,
  \href{https://arxiv.org/abs/1901.04645}{arXiv:1901.04645}.
\end{thebibliography}
\end{document}
